% Generated draft: requires mathematical and source review.
\documentclass[11pt]{article}
\usepackage[margin=25mm]{geometry}
\usepackage{fontspec}
\setmainfont{Latin Modern Roman}
\usepackage{amsmath,amssymb,mathtools,unicode-math}
\setmathfont{Latin Modern Math}
\usepackage{enumitem,needspace,xurl,hyperref}
\setlength{\parindent}{0pt}
\setlength{\parskip}{5pt}
\setlength{\emergencystretch}{3em}
\newcommand{\N}{\mathbb{N}}
\newcommand{\Z}{\mathbb{Z}}
\newcommand{\Q}{\mathbb{Q}}
\newcommand{\R}{\mathbb{R}}
\newcommand{\C}{\mathbb{C}}
\newcommand{\F}{\mathbb{F}}
\newcommand{\A}{\mathbb{A}}
\newcommand{\PP}{\mathbb{P}}
\newcommand{\E}{\mathbb{E}}
\newcommand{\eps}{\varepsilon}
\newcommand{\dd}{\,\mathrm{d}}
\begin{document}
\section*{UnsolvedMath ID: 20000399}
\textbf{Source problem code:} catalogue entry
\subsection{Definitions and mathematical statement}
A trigonal curve here means a smooth geometrically connected curve equipped with a degree-three map to \(\PP^1\), corresponding to a cubic function-field extension. Jordan Ellenberg's question asks in what senses its discriminant, under a suitable random ensemble, resembles a random polynomial. The source does not specify the field, genus/discriminant ordering, cover-versus-curve weighting or polynomial comparison measure.

The attached aid uses a discriminant-ordered function-field family in characteristic greater than three, attributed to Datskovsky--Wright/Wood. Over \(\F_q(t)\), for a finite place \(P\) of norm \(Q=q^{\deg P}\), let \(e_P\) be its field-discriminant exponent. This uses the discriminant of the maximal field extension, not an arbitrary nonmaximal cubic equation. In this tame cubic setting the values are \(0,1,2\), with exponent two corresponding to total ramification.

For every fixed finite set of finite places, the aid reports asymptotic independence and local probabilities
\[
\Pr(e_P=e)=\frac{Q^{-e}}{1+Q^{-1}+Q^{-2}},\qquad e=0,1,2.
\]
It identifies this exactly with the fixed-place law of uniformly sampled monic cubefree polynomials as their degree grows. It reports the radical-divisibility probability
\[
\Pr(P\mid\operatorname{rad}D)=\frac{Q+1}{Q^2+Q+1},
\]
which exceeds the uniform squarefree-polynomial probability \(1/(Q+1)\) by
\[
\frac{Q}{(Q+1)(Q^2+Q+1)}.
\]
On any fixed finite set, conditioning out exponent two recovers the squarefree-polynomial local law exactly:
\[
\Pr(e_P=1\mid e_P\ne2)=\frac1{Q+1}.
\]
Thus matching the full cubefree law does not mean the radical itself is uniformly squarefree.

The aid derives its global probability-generating function and finite-discriminant squarefreeness prediction
\[
\Pr(D\text{ squarefree at all finite places})\longrightarrow\frac q{q+1}
\]
only under an additional uniform-tail hypothesis. In particular, to pass from finite-place statements to exclusion of exponent two everywhere one needs a bound of the form
\[
\lim_{R\to\infty}\limsup_{\text{family size}\to\infty}
\Pr(\exists P:\deg P>R,\ e_P=2)=0.
\]
The exact family parameter, generating-function statistic and proof of this uniformity are not supplied. No corresponding assertion at infinity or in characteristics two and three is inferred.
\subsection{Short English statement}
Compare trigonal field discriminants with random polynomials. The attached aid reports an independent cubefree local law at fixed finite places, explains its difference from a uniform squarefree radical law, and keeps global squarefreeness predictions conditional on a uniform large-place tail.
\subsection{Sources}
\begin{itemize}
\item Local UnsolvedMath registry, permanent ID \(20000399\), exact public code \texttt{catalogue entry}; canonical location \texttt{aim-algebraic-number-theory-notes.json}, notes[50].
\item AIM workshop ``Arithmetic statistics over finite fields and function fields,'' Question 21, attributed to Jordan Ellenberg; supplied URL \url{https://aimath.org/pastworkshops/arithstatffieldproblems.pdf}. No year is supplied; the literature field is empty. The PDF was not consulted here.
\item Attached research attempt 1, ``A cubefree local model for trigonal discriminants,'' attributing the ensemble to Datskovsky--Wright/Wood without full publication citations.
\end{itemize}
\end{document}
